% !TEX root = ../Monografia.tex
\chapter{Metodologia}
\label{cap:metodologia}

Este capítulo descreve como a pergunta de pesquisa foi transformada em testes. A lógica é a da
Seção~\ref{sec:pergunta}: primeiro se verifica se o grafo de notícias prevê a correlação futura entre
ações (H1) e, só então, se essa informação melhora as carteiras quando entra na estimação da matriz de
covariância (H2). A Seção~\ref{sec:visao-geral} dá uma visão geral do processo, cujos detalhes de dados
e de extração ficam para os Capítulos~\ref{cap:dados} e~\ref{cap:extracao}. As seções seguintes tratam
do protocolo que protege os testes contra a seleção de resultados, das hipóteses, dos atributos
extraídos do grafo, dos testes de previsão de correlação, dos estimadores de $\Sigma$, das referências
e do placebo, das carteiras, do desenho experimental e dos testes estatísticos.

%--------------------------------------------------------------------
\section{Visão geral}
\label{sec:visao-geral}

O processo tem cinco etapas:
\begin{enumerate}
  \item \textbf{Extração.} Um LLM lê cada artigo do corpus e devolve tuplas
        \[(\text{agente}_a,\ \text{agente}_b,\ \text{tipo},\ \text{direção},\ \text{força}),\]
        datadas pela publicação do artigo (Capítulo~\ref{cap:extracao}).
  \item \textbf{Grafo.} As tuplas são filtradas, os agentes são unificados e associados a
        \emph{tickers}, e cada tupla vira um evento datado entre dois nós. O grafo numa data $t$ é a
        agregação dos eventos anteriores a $t$ (Seção~\ref{sec:atributos}).
  \item \textbf{Atributos de par.} Do grafo em $t$ saem medidas da ligação entre cada par de ações do
        universo: vínculo direto, exposição comum a entidades não negociadas e coocorrência.
  \item \textbf{Estimação de $\Sigma$.} Os atributos entram na matriz de covariância como alvo de
        \emph{shrinkage}, estrutura residual de um modelo de fatores ou visões de Black-Litterman,
        conforme o estimador (Seção~\ref{sec:estimadores}).
  \item \textbf{Carteira e avaliação.} Cada $\hat\Sigma$ alimenta um otimizador, a carteira é mantida
        até o rebalanceamento seguinte e os retornos realizados são comparados com os de estimadores de
        referência e de um placebo (Seções~\ref{sec:carteiras-metodo} a~\ref{sec:metricas}).
\end{enumerate}

A Figura~\ref{fig:pipeline} resume o fluxo. As etapas em azul são descritas no
Capítulo~\ref{cap:extracao}; as em laranja, neste capítulo.

\begin{figure}[htbp]
\centering
% Figura do pipeline (Seção 3.1). Incluída com \input; requer tikz com
% as bibliotecas arrows.meta, positioning e fit.
\begin{tikzpicture}[
  font=\footnotesize,
  node distance=6mm and 5mm,
  caixa/.style={draw, rounded corners=2pt, align=center, minimum height=13mm,
                text width=24mm, inner sep=3pt, line width=0.6pt},
  dados/.style={caixa, fill=gray!12},
  extr/.style={caixa, fill=blue!10},
  est/.style={caixa, fill=orange!14},
  seta/.style={-{Stealth[length=2.2mm]}, line width=0.7pt},
  grupo/.style={draw, dashed, rounded corners=4pt, inner sep=5pt, gray!70},
  rotulo/.style={font=\scriptsize\itshape, text=gray!80!black},
]
% linha 1: das notícias ao grafo
\node[dados] (noticias) {Notícias\\(corpus FNSPID)};
\node[extr, right=of noticias] (llm) {Extração por LLM\\(Qwen2.5-14B)};
\node[extr, right=of llm] (tuplas) {Tuplas datadas\\$(a,\,b,\,k,\,\pm,\,s)$};
\node[extr, right=of tuplas] (link) {Filtro de ruído e\\ligação a \emph{tickers}};
\node[extr, right=of link] (grafo) {Grafo dinâmico\\(eventos com\\data $<t$)};

% linha 2: do grafo à avaliação (sentido inverso)
\node[est, below=11mm of grafo] (atrib) {Atributos de par\\vínculo, exposição,\\controles};
\node[est, left=of atrib] (sigma) {Estimação\\de $\hat\Sigma$};
\node[est, left=of sigma] (cart) {Carteira\\(GMV e outras)};
\node[est, left=of cart] (aval) {Avaliação fora\\da amostra};
\node[dados, left=of aval] (ref) {Referências\\sem notícias\\e placebos};

% preços
\node[dados, below=11mm of sigma] (precos) {Preços\\(retornos diários)};

% fluxo
\draw[seta] (noticias) -- (llm);
\draw[seta] (llm) -- (tuplas);
\draw[seta] (tuplas) -- (link);
\draw[seta] (link) -- (grafo);
\draw[seta] (grafo) -- (atrib);
\draw[seta] (atrib) -- (sigma);
\draw[seta] (sigma) -- (cart);
\draw[seta] (cart) -- (aval);
\draw[seta] (ref) -- (aval);
\draw[seta] (precos) -- node[rotulo, pos=0.3, left, font=\scriptsize] {janela antes de $t$} (sigma);
\draw[seta] (precos.west) -| node[rotulo, pos=0.25, below, font=\scriptsize] {período após $t$} (aval.south);

% grupos e capítulos
\node[grupo, fit=(llm)(grafo)] (g1) {};
\node[rotulo, anchor=south east] at (g1.north east) {Capítulo~\ref{cap:extracao}};
\node[grupo, fit=(atrib)(aval)] (g2) {};
\node[rotulo, anchor=north east] at (g2.south east) {Capítulo~\ref{cap:metodologia}};
\end{tikzpicture}

\caption{Visão geral do processo, das notícias à avaliação das carteiras. Em cinza, as entradas; em
azul, a extração e a construção do grafo; em laranja, a estimação e a avaliação. Os preços entram na
estimação apenas com a janela anterior a $t$ e, na avaliação, com o período posterior.}
\label{fig:pipeline}
\end{figure}

Uma regra vale para todas as etapas: nada do que entra numa decisão tomada em $t$ pode depender de
informação posterior a $t$. Como a data de publicação das notícias do corpus não tem hora
(Capítulo~\ref{cap:dados}), uma notícia do dia $d$ pode ter saído depois do fechamento; o grafo em $t$
usa, por isso, apenas eventos com data estritamente anterior a $t$.

%--------------------------------------------------------------------
\section{Protocolo de pré-registro}
\label{sec:preregistro}

Comparar vários estimadores, carteiras e parâmetros no mesmo período cria o risco de encontrar, por
acaso, uma combinação que parece funcionar. Para controlar esse risco, o trabalho adotou um protocolo de
pré-registro incremental:

\begin{itemize}
  \item cada desenho de teste foi escrito e versionado no repositório do projeto \emph{antes} de ser
        executado com os dados reais, com o modelo, as referências, a estatística, o nível de
        significância e o critério de sucesso fixados;
  \item qualquer mudança feita depois da execução foi registrada como tal, com o motivo, e os
        resultados das versões anteriores foram mantidos;
  \item quando um desenho falhou e um novo desenho foi proposto para a mesma hipótese, o nível de
        significância do novo foi reduzido pela correção de Bonferroni ($0{,}05/k$ para o $k$-ésimo
        desenho), de modo que tentar de novo não aumenta a chance de um falso positivo;
  \item os parâmetros de construção do grafo (janela de 365 dias e meia-vida de 90 dias) e a janela de
        estimação (252 pregões) foram fixados antes do primeiro teste e não foram calibrados;
  \item todo resultado de carteira é confrontado com um placebo e confirmado em universos de outros
        tamanhos, e os primeiros desenhos foram refeitos depois que a amostra de artigos foi ampliada.
\end{itemize}

Os desvios em relação aos pré-registros são declarados no ponto do texto em que ocorrem, e as análises
feitas depois de conhecido um resultado são identificadas como exploratórias.

%--------------------------------------------------------------------
\section{Hipóteses}
\label{sec:hipoteses}

As duas hipóteses gerais foram apresentadas na Seção~\ref{sec:pergunta}. H1 diz que o grafo contém
informação sobre a correlação futura; H2, que essa informação melhora as carteiras. Para que H1 possa ser
testada, é preciso dizer \emph{por qual canal} o grafo carregaria essa informação. A teoria de finanças
oferece dois canais, que dão as duas hipóteses específicas de H1. São, ao todo, três hipóteses,
resumidas na Tabela~\ref{tab:hipoteses}.

\paragraph{H1a: exposição comum.} Nos modelos de fatores, duas ações covariam porque estão expostas às
mesmas fontes de risco sistemático \citep{ross1976,chan1999}. Parte dessas fontes são entidades que não
são negociadas como ações, como a política monetária, o preço do petróleo ou a economia chinesa, e as
notícias registram quais empresas são afetadas por cada uma. Se duas empresas aparecem ligadas às mesmas
entidades, elas devem ter cargas parecidas nos fatores correspondentes e, portanto, correlação maior. A
exposição comum $g_{ij}$ da Equação~\eqref{eq:exposicao} mede essa semelhança, e a hipótese é que ela
preveja a correlação futura do par além da correlação histórica e do setor. O setor de
\citet{elton1973} é o controle natural, porque também é uma forma de exposição comum, só que fixa e
grosseira.

\paragraph{H1b: tipo de relação.} A literatura de vínculos econômicos mostra que relações diretas entre
empresas transmitem choques, e que o canal depende do tipo de relação. Clientes e fornecedores dividem
o mesmo fluxo de demanda \citep{cohen2008,menzly2010}; concorrentes disputam o mesmo mercado de produtos e
respondem aos mesmos choques de demanda e de custo \citep{hoberg2016}; empresas com tecnologias próximas
compartilham choques de inovação \citep{lee2019}; e participações acionárias e aquisições ligam
diretamente os fluxos de caixa. Nem todo vínculo, porém, deveria ter o mesmo efeito: uma disputa
judicial ou uma ação regulatória não cria exposição econômica duradoura. A hipótese é que os efeitos dos
grupos de relação sobre a correlação residual futura sejam diferentes entre si. Como cada relação
extraída também tem um sinal, a análise de H1b examina ainda se o sinal atribuído pelo LLM indica a
direção do co-movimento.

\paragraph{H2: as notícias como \emph{prior} da correlação.} Na leitura bayesiana do \emph{shrinkage}, o
alvo é a média de um \emph{prior} sobre $\Sigma$, e um alvo mais próximo da covariância verdadeira dá uma
estimativa melhor \citep{ledoit2004honey}. \citet{becquin2023} mostram que um alvo construído com texto
pode ajudar, e \citet{lu2023}, que uma rede externa aos retornos pode estruturar a parte da covariância que
o mercado não explica. Se H1 vale, as notícias preveem a correlação de cada par, e essa previsão é um
\emph{prior} candidato. A hipótese é que a carteira de variância mínima construída com o \emph{prior} de
setor e notícias tenha variância realizada menor que a construída com o \emph{prior} só de setor, a
referência que encarna \citet{elton1973}, e que o ganho venha do conteúdo das notícias: a carteira deve
superar também um placebo em que os atributos de notícia são embaralhados entre as ações. H2 foi testada
numa sequência de nove desenhos (Tabela~\ref{tab:desenhos}), que variam o modo de levar as notícias à
matriz e o otimizador.

\begin{table}[htbp]
\centering
\small
\caption{Hipóteses do trabalho, fundamentação e teste.}
\label{tab:hipoteses}
\begin{tabularx}{\textwidth}{>{\raggedright\arraybackslash}p{0.9cm}>{\raggedright\arraybackslash}X>{\raggedright\arraybackslash}p{3.6cm}>{\raggedright\arraybackslash}p{3.6cm}}
\toprule
& hipótese & fundamentação & teste \\
\midrule
H1a & a exposição comum a entidades não negociadas prevê a correlação futura do par além do histórico e do setor & modelos de fatores; \citet{elton1973}, \citet{chan1999} & ganho fora da amostra (escada E0--E4) e coeficiente de $g_{ij}$ nos testes de pares \\
H1b & os tipos de relação direta têm efeitos diferentes sobre a correlação residual futura & \citet{cohen2008}, \citet{menzly2010}, \citet{hoberg2016} & Wald da igualdade dos coeficientes de grupo \\
H2 & a variância mínima com o \emph{prior} de setor e notícias tem variância realizada menor que com o \emph{prior} de setor e que com o placebo & \citet{ledoit2004honey}, \citet{becquin2023}, \citet{lu2023} & razão de variâncias, placebo e confirmação em outros universos \\
\bottomrule
\end{tabularx}
\end{table}

O trabalho se limita a essas três hipóteses de propósito. Cada uma tem um mecanismo econômico descrito na
literatura, uma previsão fixada antes do teste e uma estatística própria, e cada uma é examinada a fundo:
com controles que separam o canal das explicações mais simples (histórico, setor, atenção da mídia), com
inferência adequada à dependência entre pares e, nas carteiras, com placebo e confirmação em universos de
tamanhos diferentes. Investigar poucas hipóteses bem fundamentadas
com esse cuidado informa mais do que testar muitas ideias com menos rigor, e reduz o risco de que um
resultado positivo surja apenas do número de tentativas.

%--------------------------------------------------------------------
\section{Atributos de par a partir do grafo}
\label{sec:atributos}

\subsection{O grafo numa data}

Cada relação extraída e filtrada é um evento $e$ entre dois nós, com data $t_e$, classe de relação
$c_e$, sinal $\mathrm{sgn}_e\in\{-1,+1\}$ e força $s_e\in[0,1]$. O grafo na data $t$ agrega os eventos
da janela $[t-W,\,t)$, com $W=365$ dias corridos, ponderados por um decaimento exponencial de meia-vida
$\tau$:
\begin{equation}
  w_e(t) = \mathrm{sgn}_e\, s_e\, 0{,}5^{\,(t-t_e)/\tau}.
  \label{eq:peso-evento}
\end{equation}
A meia-vida é $\tau=90$ dias em todos os testes, valor fixado sem calibração.

\subsection{Vínculo direto}

Para duas ações $i$ e $j$ do universo, o vínculo direto soma os eventos entre elas, nos dois sentidos.
Na forma com sinal, $w_{ij}=\sum_e w_e(t)$; separando por sinal, $w^+_{ij}$ e $w^-_{ij}$ são as somas,
em módulo, dos eventos positivos e negativos; e, por grupo $k$ de classes de relação,
$d^{(k)}_{ij}=\sum_{e:\,c_e\in k}|w_e(t)|$. Os grupos reúnem classes com efeito econômico parecido sobre
o co-movimento. Nas regressões que alimentam as carteiras, são cinco: cadeia (\emph{supply chain} e
parceria), concorrência, controle (participação acionária e fusões e aquisições), impacto e outros. Nos
testes de H1b, a cadeia e a parceria são separadas e há um grupo legal e regulatório, o que dá seis
grupos mais o residual. Como os vínculos têm cauda longa, entram nas regressões como $\log(1+d)$ ou, com
sinal, como $\mathrm{sgnlog}(w)=\mathrm{sgn}(w)\log(1+|w|)$.

\subsection{Exposição comum}

Arestas diretas entre ações são raras, mas muitos pares compartilham alguma entidade que não é
negociada, como o Federal Reserve, o petróleo ou a China (Capítulo~\ref{cap:extracao}). A exposição
comum mede essa ligação indireta. Seja $\tilde A$ a matriz ações $\times$ entidades fora do universo, com
$\tilde A_{im}$ igual à soma dos eventos entre a ação $i$ e a entidade $m$. Cada coluna é ponderada pelo
inverso da frequência da entidade (IDF),
\[
  \mathrm{idf}_m=\log\frac{n_t}{\mathrm{df}_m},
\]
em que $n_t$ é o número de empresas com \emph{ticker} presentes no grafo da janela e $\mathrm{df}_m$ é o
número delas ligadas a $m$. Assim, uma entidade ligada a quase todas as empresas, como ``S\&P 500'',
pesa pouco, e ``petróleo'' pesa para as petroleiras. Com as linhas normalizadas para norma unitária, a
exposição comum é o cosseno
\begin{equation}
  g_{ij} = \bigl(\tilde A\tilde A^\top\bigr)_{ij} = \cos(\tilde A_i,\tilde A_j).
  \label{eq:exposicao}
\end{equation}
A matriz $G=[g_{ij}]$ é a projeção do grafo bipartido ações--entidades sobre as ações. Uma ação sem
nenhuma entidade na janela tem $g_{ij}=0$ com todas as outras. Nas regressões do universo das ações mais
citadas, $\tilde A$ usa o módulo dos eventos, e $g_{ij}\in[0,1]$; nos testes de pares do universo do
S\&P~500, usa o sinal, e $g_{ij}\in[-1,1]$, negativo quando as duas ações se ligam às mesmas entidades com
efeitos opostos.

\subsection{Controles}

Três atributos separam o conteúdo das relações de efeitos mais simples:
\begin{itemize}
  \item \textbf{coocorrência}: número de artigos extraídos em que $i$ e $j$ aparecem em algum evento,
        com ou sem relação entre si; separa ``o LLM encontrou uma relação'' de ``saíram no mesmo
        texto'';
  \item \textbf{co-menção}: número de artigos do corpus marcados com os dois \emph{tickers}, a variável
        de \citet{tang2019};
  \item \textbf{setor}: indicadores de mesmo setor na classificação Fama-French de 49 indústrias e no
        código SIC de dois dígitos, o controle de \citet{elton1973}.
\end{itemize}
Os testes de pares acrescentam a atenção da mídia, $\log(1+m_i)+\log(1+m_j)$ e
$|\log(1+m_i)-\log(1+m_j)|$, em que $m_i$ é o número de eventos da ação $i$ na janela, e a distância de
tamanho entre as duas empresas.

%--------------------------------------------------------------------
\section{Previsão de correlação e testes de pares}
\label{sec:metodo-h1}

H1 é testada de dois modos complementares. O primeiro mede, fora da amostra, quanto os atributos de
notícia melhoram a previsão da correlação de cada par, no universo das ações mais citadas. O segundo
estima o efeito de cada atributo com inferência adequada a dados de pares, no universo do S\&P~500, que
não depende das notícias (Capítulo~\ref{cap:dados}).

\subsection{Regressão preditiva de correlação}

Em cada data de rebalanceamento $t$, para cada par $(i,j)$ do universo, a variável dependente é a
correlação realizada no período de carteira seguinte, $\rho^{\mathrm{fut}}_{ij}$, na escala de Fisher
$z(\rho)=\operatorname{atanh}(\rho)$, que estabiliza a variância. A especificação mais completa é
\begin{equation}
  z(\rho^{\mathrm{fut}}_{ij}) = a + b\,z(\rho^{\mathrm{hist}}_{ij})
  + \sum_k \beta_k\log\bigl(1+d^{(k)}_{ij}\bigr) + \gamma\, g_{ij}
  + \theta^\top x_{ij} + \varepsilon_{ij},
  \label{eq:h1}
\end{equation}
em que $\rho^{\mathrm{hist}}_{ij}$ é a correlação nos 252 pregões anteriores a $t$ e $x_{ij}$ reúne
coocorrência, co-menção e setor. Os coeficientes $\beta_k$ e $\gamma$ traduzem cada tipo de relação em
correlação. As especificações testadas formam uma escada:
\begin{itemize}
  \item E0, persistência: só $z(\rho^{\mathrm{hist}})$;
  \item E1: E0 com co-menção;
  \item E2: E0 com coocorrência;
  \item E3: E0 com vínculos diretos por grupo e exposição comum;
  \item E4: todos os atributos.
\end{itemize}
E0, E1 e E4 também são estimadas com os indicadores de setor. A comparação central é entre E0 com setor e
E4 com setor: a diferença mede o que as notícias acrescentam ao histórico e ao setor.

O desempenho é medido fora da amostra. Em cada $t$, os coeficientes são estimados no painel empilhado
dos rebalanceamentos anteriores (janela expansiva, com ao menos 24 meses) e aplicados aos pares de $t$.
A variável dependente e os regressores são centrados em cada mês, de modo que a regressão explica a
ordenação dos pares dentro do mês, e não o nível médio da correlação. O $R^2$ fora da amostra
\citep{campbell2008} é
\[
  R^2_{oos}=1-\frac{\sum_t\sum_{(i,j)}(y_{ij,t}-\hat y_{ij,t})^2}{\sum_t\sum_{(i,j)}(y_{ij,t}-\bar y_t)^2},
\]
com as somas sobre todos os pares e meses de avaliação. O ganho de uma especificação sobre outra é a
diferença dos seus $R^2_{oos}$, testada com erro-padrão de \citet{newey1987} na série mensal das
diferenças de erro quadrático.

\subsection{Correlação residual}

A correlação total entre ações é dominada pelo mercado. Para isolar o que é específico de cada par,
os testes também usam a correlação residual a um fator de mercado. Seja $m_t$ a média igual dos
retornos do universo no dia. Na janela de estimação, cada ação é regredida no mercado,
\begin{equation}
  x_{it} = a_i + b_i\,m_t + e_{it},
  \label{eq:fator-mercado}
\end{equation}
e a correlação residual realizada no período seguinte é calculada com os resíduos fora da amostra, com
os $b_i$ da janela e o $m_t$ observado. A escada E0--E4 é estimada para a correlação total e para a
residual.

\subsection{Testes de pares}
\label{sec:testes-pares}

Os testes de pares usam a correlação residual dos $h=21$ pregões seguintes como variável dependente e
todos os meses de 2011 a 2023. A regressão é
\begin{equation}
  z(\rho^{\mathrm{fut}}_{e,ij}) = a + b\,z(\rho^{\mathrm{hist}}_{e,ij}) + \gamma\,g_{ij}
  + \sum_k \beta_k\log\bigl(1+d^{(k)}_{ij}\bigr) + \theta^\top c_{ij} + \varepsilon_{ij},
  \label{eq:pares}
\end{equation}
em que $\rho_e$ é a correlação residual e $c_{ij}$ reúne os indicadores de mesmo setor (Fama-French~49 e
SIC de dois dígitos), a distância de tamanho e os dois termos de atenção. Na análise do sinal, a soma
sobre os grupos é trocada por $\beta\,\mathrm{sgnlog}(w_{ij})$ ou por
$\beta^+\log(1+w^+_{ij})+\beta^-\log(1+w^-_{ij})$.

Pares que compartilham uma ação não são independentes. Por isso a inferência usa a regressão de
\citet{fama1973}: uma regressão em corte transversal em cada mês, com os regressores padronizados no
mês, e a média dos coeficientes mensais, com erro-padrão de \citet{newey1987} com $L=3+h/21$ defasagens.
Cada coeficiente mede, assim, a variação de $z(\rho)$ por um desvio-padrão do regressor. Como segunda
inferência, o procedimento QAP \citep{krackhardt1988} permuta juntas as linhas e colunas de todas as
matrizes de notícia em cada mês e refaz o Fama-MacBeth; o $p$ é a fração de permutações com estatística ao
menos tão extrema quanto a observada, com 1.000 permutações, o que limita o menor $p$ possível a
$1/1001\approx0{,}001$. O $p$ de cada teste é o maior entre o de Fama-MacBeth e o do QAP. A
Tabela~\ref{tab:testes-diadicos} resume a estatística de cada teste.

\begin{table}[htbp]
\centering
\small
\caption{Estatísticas e critérios dos testes de pares.}
\label{tab:testes-diadicos}
\begin{tabularx}{\textwidth}{>{\raggedright\arraybackslash}p{1.6cm}>{\raggedright\arraybackslash}X>{\raggedright\arraybackslash}p{4.2cm}}
\toprule
& estatística & sucesso \\
\midrule
H1a & coeficiente $\gamma$ da exposição comum, com a aresta direta e os controles & positivo e significativo \\
H1b & teste de Wald da igualdade dos seis coeficientes de grupo $\beta_k$, com covariância de Newey-West & igualdade rejeitada \\
H1b, sinal & coeficiente de $\mathrm{sgnlog}(w)$ e diferença $\beta^+-\beta^-$ & positivos e significativos \\
\bottomrule
\end{tabularx}
\end{table}

%--------------------------------------------------------------------
\section{Estimadores da matriz de covariância com notícias}
\label{sec:estimadores}

Três modos de levar as notícias à matriz de covariância foram testados. Em todos, as variâncias vêm dos
preços, e as notícias atuam apenas sobre as correlações.

\subsection{Alvo de shrinkage com notícias}

O primeiro modo segue \citet{becquin2023}: as notícias definem a correlação-alvo. A correlação prevista
para o par $(i,j)$ é
\begin{equation}
  R^T_{ij} = \tanh\!\bigl(\bar z + \beta^\top f_{ij}\bigr),\qquad R^T_{ii}=1,
  \label{eq:alvo-prior}
\end{equation}
em que $\bar z$ é a correlação média da janela em escala $z$, $f_{ij}$ são os atributos de notícias (e de
setor) e $\beta$ é estimado por mínimos quadrados no painel dos meses anteriores, com $y$ e $f$
centrados no mês. A correlação histórica do par não entra em \eqref{eq:alvo-prior}: com ela, o alvo
dependeria da mesma amostra que $S$, o que viola a hipótese de alvo fixo da fórmula de $\delta$. Com
$\beta=0$, o alvo é a correlação constante de \citet{ledoit2004honey}; como $\bar z$ é a média da janela,
o alvo preserva o nível médio da correlação, isto é, o modo de mercado. Se $R^T$ não for positiva
semidefinida, é substituída pela correlação mais próxima \citep{higham2002}. A covariância-alvo é
$T=D^{1/2}R^TD^{1/2}$, com $D=\operatorname{diag}(S)$, e o estimador é a Equação~\eqref{eq:shrinkage}.

O $\delta$ ótimo generaliza a fórmula de \citet{ledoit2004honey}, trocando a correlação constante
$\bar r$ pela correlação-alvo $r^T_{ij}$ de cada par:
\begin{align}
  \hat\pi &= \sum_{i,j}\widehat{\mathrm{AsyVar}}\bigl(\sqrt N s_{ij}\bigr), \qquad
  \hat\gamma = \|T-S\|_F^2, \nonumber\\
  \hat\rho &= \sum_i\hat\pi_{ii} + \sum_{i\neq j}\frac{r^T_{ij}}{2}
     \left(\sqrt{\frac{s_{jj}}{s_{ii}}}\,\hat\vartheta_{ii,ij} + \sqrt{\frac{s_{ii}}{s_{jj}}}\,\hat\vartheta_{jj,ij}\right),
  \label{eq:delta-geral}\\
  \delta^* &= \max\!\left\{0,\ \min\!\left\{\frac{\hat\pi-\hat\rho}{N\hat\gamma},\ 1\right\}\right\}, \nonumber
\end{align}
em que $\hat\pi_{ij}=\frac1N\sum_t[(r_{it}-\bar r_i)(r_{jt}-\bar r_j)-s_{ij}]^2$ estima a variância
assintótica de $\sqrt N s_{ij}$ e
$\hat\vartheta_{ii,ij}=\frac1N\sum_t[(r_{it}-\bar r_i)^2-s_{ii}][(r_{it}-\bar r_i)(r_{jt}-\bar r_j)-s_{ij}]$
estima a covariância assintótica entre $\sqrt N s_{ii}$ e $\sqrt N s_{ij}$. Como em
\citet{ledoit2004honey}, o alvo usa as variâncias amostrais, as covariâncias entre as variâncias e os
elementos fora da diagonal entram em $\hat\rho$, e $S$ é calculada com divisor $N$ nessas fórmulas.

\subsection{Fator de mercado e correlação residual}

O segundo modo aplica as notícias à parte da covariância que o mercado não explica, como na
Equação~\eqref{eq:fatores} com um fator. Com o modelo \eqref{eq:fator-mercado} estimado na janela,
\begin{equation}
  \hat\Sigma = \hat\sigma_m^2\, b\, b^\top + \hat\Sigma_e,
  \label{eq:fator-residuo}
\end{equation}
em que $\hat\sigma_m^2$ é a variância do mercado, $b$ é o vetor de betas e $\hat\Sigma_e$ é a covariância
dos resíduos $E$, encolhida para uma correlação residual \emph{a priori} $R_{\mathrm{prior}}$ pela fórmula
\eqref{eq:delta-geral}. O que muda entre os estimadores é $R_{\mathrm{prior}}$:
\begin{itemize}
  \item identidade: o modelo de fator único puro;
  \item correlação constante $\bar r_e$: a referência estatística;
  \item $\tanh(\bar z_e+\beta^\top\text{setor})$: a \textbf{referência a bater};
  \item $\tanh(\bar z_e+\beta^\top[\text{setor},\text{notícias}])$: o \textbf{modelo principal};
  \item o anterior com os atributos de notícia permutados entre as ações: o placebo.
\end{itemize}
Aqui, $\beta$ é estimado com $y$ igual à correlação residual realizada no mês. É o desenho de
\citet{lu2023}, com a rede de notícias no lugar da rede de co-negociação. Como o fator de mercado fica
fora do \emph{prior}, as notícias não precisam descrever o modo comum a todas as ações, apenas o que os
pares têm em comum além dele, que é onde H1 mostra que elas informam.

\subsection{Black-Litterman na covariância}

Nos dois modos anteriores, um único $\delta$ vale para todos os pares: um par com centenas de eventos e
um par ligado só por exposição comum recebem o mesmo peso da notícia. O terceiro modo dá a cada par a
confiança que a sua evidência justifica, na estrutura de \citet{black1992} aplicada à covariância
\citep{qian2001}. Apenas a covariância residual $\Sigma_e$ da Equação~\eqref{eq:fator-residuo} é
atualizada.

\begin{itemize}
  \item \textbf{Prior.} $\Sigma_0$ é a covariância residual do modelo com setor, a referência a bater,
        com variâncias $s_i=(\Sigma_0)_{ii}$. A incerteza do prior sobre os elementos fora da diagonal é
        a covariância de Wishart da covariância amostral \citep{anderson2003}, avaliada em $\Sigma_0$
        com $N=252$ pregões:
        \begin{equation}
          V_{(ij),(kl)} = \frac{\sigma_{0,ik}\,\sigma_{0,jl} + \sigma_{0,il}\,\sigma_{0,jk}}{N},
          \label{eq:wishart}
        \end{equation}
        o análogo, na covariância, da convenção de \citet{he1999} de tomar a incerteza do \emph{prior}
        como o inverso do tamanho da amostra. Uma visão sobre o par $(i,j)$ move também os pares $(i,l)$
        e $(j,l)$, na proporção das covariâncias do prior.
  \item \textbf{Visões.} Há uma visão para cada par com algum atributo de notícia não nulo. Seu valor é a
        correlação prevista pelo modelo com setor e notícias, na escala da covariância:
        $q_{ij}=\tanh(\hat z_{ij})\sqrt{s_is_j}$. A matriz $P$ seleciona o elemento $(i,j)$ de
        $\operatorname{vec}\Sigma_e$, de modo que a visão é linear em $\Sigma_e$.
  \item \textbf{Confiança.} Cada par com visão cai numa classe de evidência $c$: só exposição comum ou
        co-menção ($c=0$), ou um dos quatro quartis de coocorrência ($c=1,\dots,4$). Para cada classe,
        o erro da visão em escala $z$ é estimado no painel dos meses de treino, descontado o ruído
        amostral de uma correlação de $d_m$ pregões,
        \begin{equation}
          \omega_c = \overline{\hat e^2}\big|_c - \overline{1/(d_m-3)},\qquad \omega_c\ge \frac{1}{N-3},
          \label{eq:omega}
        \end{equation}
        e levado à escala da covariância pelo método delta,
        $\Omega_{ij}=(1-\tanh^2\hat z_{ij})^2\,s_is_j\,\omega_{c(ij)}$.
  \item \textbf{Posterior.} É a média do Black-Litterman na forma de filtro de Kalman,
        \begin{equation}
          \operatorname{vec}\Sigma_e^{BL} = \operatorname{vec}\Sigma_0
          + VP^\top\bigl(PVP^\top+\Omega\bigr)^{-1}\bigl(q-P\operatorname{vec}\Sigma_0\bigr),
          \label{eq:bl}
        \end{equation}
        resolvida por gradiente conjugado sem formar $V$, e projetada para a correlação positiva
        semidefinida mais próxima, com as variâncias de volta.
\end{itemize}

%--------------------------------------------------------------------
\section{Estimadores de referência e placebo}
\label{sec:referencias}

Cada estimador com notícias tem uma referência sem notícias que só difere dele pela informação do
grafo: o alvo de setor, no primeiro modo, e o fator de mercado com resíduo encolhido para o setor, nos
outros dois. Essa referência encarna o resultado de \citet{elton1973} e é a hipótese nula natural para um
grafo que, como se verá, liga sobretudo empresas do mesmo setor. Para situar os resultados na literatura,
as tabelas trazem também as referências estatísticas usuais: a carteira $1/N$ \citep{demiguel2009}, a
covariância amostral, o \emph{shrinkage} linear de Ledoit e Wolf com os três alvos clássicos (identidade,
correlação constante e fator único) e o \emph{shrinkage} não linear \citep{ledoit2020}.

Superar a referência não basta: o ganho poderia vir apenas de regularizar a matriz com qualquer
estrutura esparsa. O placebo testa se o \emph{conteúdo} das relações importa. Em cada mês, as ações são
permutadas ao acaso nas matrizes de atributos de notícia, com a mesma permutação em todos os atributos,
enquanto o setor e os retornos ficam no lugar; os coeficientes $\beta$ são os treinados com os dados
reais. O placebo tem, portanto, a mesma distribuição de vínculos, a mesma esparsidade e a mesma
intensidade de \emph{shrinkage} que as notícias verdadeiras, mas liga as empresas erradas. Se as notícias
fossem só regularização genérica, o placebo empataria com elas.

%--------------------------------------------------------------------
\section{Carteiras}
\label{sec:carteiras-metodo}

Os estimadores foram avaliados nas seguintes carteiras, todas com soma dos pesos igual a um:
\begin{itemize}
  \item \textbf{GMV sem restrição}, Equação~\eqref{eq:gmv}: a carteira em que o erro de $\hat\Sigma$
        passa inteiro para os pesos; é a carteira principal;
  \item \textbf{variância mínima restrita}: sem venda a descoberto e com teto de 10\% por ativo,
        resolvida por gradiente projetado (FISTA);
  \item \textbf{Markowitz}, Equação~\eqref{eq:markowitz} com $\gamma=5$, e \textbf{máximo Sharpe}, com
        retornos esperados de Bayes-Stein \citep{jorion1986} calculados na janela;
  \item \textbf{máxima diversificação} \citep{choueifaty2008};
  \item \textbf{HRP} \citep{lopezdeprado2016}, com a distância de correlação $\sqrt{(1-\rho_{ij})/2}$, a
        mesma da Seção~\ref{sec:redes} a menos de uma constante;
  \item \textbf{mínimo \emph{tracking error}} contra o SPY, com a covariância entre ativos e índice
        estimada na janela.
\end{itemize}
Os problemas de razão (máximo Sharpe e máxima diversificação) com restrições foram resolvidos pela
reformulação convexa homogênea \citep{cornuejols2007} por pontos interiores.

%--------------------------------------------------------------------
\section{Desenho experimental}
\label{sec:desenho}

\subsection{Configuração comum}

Todas as avaliações de carteira seguem a mesma configuração:
\begin{itemize}
  \item universo das $n$ ações mais citadas nas notícias (Capítulo~\ref{cap:dados}), com $n=100$ no
        teste principal e $n=50$ e $n=200$ como confirmação; com a janela de 252 pregões, há sempre mais
        observações que ações;
  \item rebalanceamento no primeiro pregão de cada mês, com a carteira mantida até o rebalanceamento
        seguinte;
  \item janela de estimação de 252 pregões anteriores ao rebalanceamento;
  \item período avaliado de janeiro de 2013 a dezembro de 2023; os anos de 2011 e 2012 servem apenas
        para treinar os coeficientes que transformam notícias em correlação, que exigem ao menos 24
        meses de histórico;
  \item coeficientes treinados só com meses cujo período de carteira terminou antes de $t$;
  \item custo de transação de 10 pontos-base sobre o giro, medido contra os pesos que a carteira teria
        depois de variar com os preços desde o último rebalanceamento;
  \item SPY como índice de referência.
\end{itemize}

\subsection{Os desenhos testados para H2}

A Tabela~\ref{tab:desenhos} lista os desenhos na ordem em que foram testados, com o critério de sucesso
de cada um. Cada desenho surgiu do resultado do anterior, e o nível de significância caiu a cada nova
tentativa (Seção~\ref{sec:preregistro}). O D1 foi exploratório; o D2 é o teste principal de H2, com três
critérios fixados antes da execução; os demais perguntam se o resultado do D2 muda com restrições de
peso, com outro otimizador ou com outro modo de combinar as notícias.

\begin{table}[htbp]
\centering
\small
\caption{Desenhos testados para H2, na ordem de execução.}
\label{tab:desenhos}
\begin{tabularx}{\textwidth}{>{\raggedright\arraybackslash}p{0.9cm}>{\raggedright\arraybackslash}p{3.6cm}X}
\toprule
& desenho & critério de sucesso \\
\midrule
D1 & alvo de notícias, GMV & exploratório, sem critério fixado \\
D2 & fator + resíduo, GMV sem restrição & (i) as notícias preveem a correlação residual além do setor, com $t>2$; (ii) razão de variâncias contra o setor $<1$ com $p<0{,}05$ e abaixo do placebo em $n=100$; (iii) razão $<1$ em $n=50$ e $200$ e com a amostra ampliada de notícias \\
D3 & fator + resíduo, variância mínima sem venda a descoberto e teto de 10\% & razão $<1$ com $p<0{,}025$ e abaixo do placebo em $n=100$; razão $<1$ em $n=50$ e $200$ \\
D4--D8 & fator + resíduo em Markowitz, máximo Sharpe, máxima diversificação, HRP e mínimo \emph{tracking error} & melhora com $p<0{,}01$ por otimizador e melhor que o placebo; mesma direção em $n=50$ e $200$ \\
D9 & Black-Litterman na covariância, GMV sem restrição & razão $<1$ com $p<0{,}0125$ e abaixo do placebo em $n=100$; razão $<1$ em $n=50$ e $200$ \\
\bottomrule
\end{tabularx}
\end{table}

A melhora é medida pela razão de variâncias dos retornos líquidos (variância mínima, máxima
diversificação, HRP), pela diferença de Sharpe líquido (Markowitz, máximo Sharpe) ou pela razão de
variâncias do retorno ativo (mínimo \emph{tracking error}).

%--------------------------------------------------------------------
\section{Métricas e testes estatísticos}
\label{sec:metricas}

\paragraph{Razão de variâncias.} A estatística principal das carteiras é a razão entre as variâncias dos
retornos diários líquidos de duas carteiras nos pregões em comum, com a carteira com notícias no numerador:
valores abaixo de 1 favorecem as notícias. O $p$-valor bicaudal e o intervalo de confiança vêm de um
\emph{bootstrap} circular em blocos de 10 dias, com 2.000 reamostragens dos pares de retornos
\citep{ledoit2011}, o que preserva a dependência temporal e a correlação entre as duas carteiras. A
diferença de índices de Sharpe é testada da mesma forma \citep{ledoit2008}.

\paragraph{Erro de correlação.} A qualidade de $\hat\Sigma$ como previsão é medida pelo erro quadrático
médio entre as correlações estimadas e as realizadas no mês seguinte. A diferença entre dois estimadores
é testada na série mensal, com erro-padrão de Newey-West, como no teste de \citet{diebold1995}.

\paragraph{Critérios e tentativas repetidas.} H1a tem dois critérios: ganho fora da amostra positivo com
$t>2$ na escada E0--E4 e coeficiente positivo e significativo a 5\% nos testes de pares. H1b é avaliada
pelo teste de Wald, e os coeficientes individuais dos grupos são corrigidos pelo método de
\citet{benjamini1995}. Para H2, o nível de significância cai a cada novo desenho pela correção de
Bonferroni (Seção~\ref{sec:preregistro}).

\paragraph{Diagnóstico de estabilidade.} O número de condição $\lambda_{\max}/\lambda_{\min}$ da correlação
estimada é registrado mês a mês para cada estimador \citep{lopezdeprado2016}. Ele não é critério de
sucesso, mas mede quanto a inversão da matriz amplifica o erro e ajuda a explicar por que um estimador
que prevê melhor as correlações pode produzir uma carteira pior.

\clearpage
